\newcommand{\letitle}{Stage L3 : Formalisation de théorie des jeux en Coq}
\newcommand{\leauthor}{Ulysse Durand\\encadré par\\Etienne Miquey\\et\\Pierre Clairambault}

\documentclass{beamer}
\usepackage{tikzit}

\setbeamertemplate{footline}[frame number]
\usepackage[utf8]{inputenc}
\usepackage{amssymb}
\usepackage{amsmath}
\usepackage{xcolor}
\usepackage{enumitem}
\usepackage{bbold}
\usepackage[xcolor,leftbars]{changebar}
\usepackage{ebproof}
\setlength{\parindent}{0pt}
\setcounter{secnumdepth}{1}
\newenvironment{myindentpar}
 {\begin{list}{}
         {\setlength{\leftmargin}{1em}}
         \item[]
 }
 { \end{list}}

\definecolor{DarkBlue}{RGB}{0,16,80}
\newcommand{\norm}[1]{\lvert #1 \rvert}
\newsavebox{\mybox}
\newlength{\mydepth}
\newlength{\myheight}
\newenvironment{answer}
{\par\begin{lrbox}{\mybox}\quad\begin{minipage}{\linewidth}\color{black}\setlength{\parskip}{10pt plus 1pt minus 1pt}\vspace*{-.7\baselineskip}}
{\end{minipage}\end{lrbox}
\settodepth{\mydepth}{\usebox{\mybox}}
\settoheight{\myheight}{\usebox{\mybox}}
\addtolength{\myheight}{\mydepth}
\noindent\makebox[0pt]{
  \color{gray}\hspace{-0pt}\rule[-\mydepth]{1pt}{\myheight}}
  \usebox{\mybox}
  }

\usepackage{listings}
\lstset
{
    language=Caml,
    basicstyle=\footnotesize,
    numbers=left,
    stepnumber=1,
    showstringspaces=false,
    tabsize=1,
    breaklines=true,
    breakatwhitespace=false,
}

\usepackage{hyperref}
\setlength{\parskip}{0.15cm}
\hypersetup{
    colorlinks=truem,
    linkcolor=black,
    filecolor=red,
    urlcolor=blue
}

\urlstyle{same}

\everymath{\displaystyle}
\title{\letitle}
\author{\leauthor}
\institute{}
\date{}
\begin{document}

\maketitle
\begin{frame}
\frametitle{Sommaire}

Théorie des jeux, utile pour la théorie de la programmation

Nous allons la formaliser (catégorie)
\pause
\begin{itemize}
\item Définitions ensemblistes
\item Définitions Coq (apport perso)
\item Conclusion
\end{itemize}
\end{frame}


\begin{frame}
\frametitle{Structure d'événement (S.E.) :}

$E = (|E|, \leq_E, \#_E)$
\begin{itemize}
\item $\leq_E$ (la causalité) : ordre partiel sur $E$
\item $\#_E$ (le conflit) : relation binaire symétrique irréflexive
\begin{itemize}
\item à causes finies
\item qui vérifie l'axiome de vendetta
\end{itemize}
\end{itemize}
\pause

\tikzfig{ex1}

$\sim$ étant la relation de conflit immédiat et\\
$\rightarrow$ la relation de causalité immédiate
\end{frame}

\begin{frame}
\frametitle{Configuration :}


$x$ fini $\subset |E|$
\begin{itemize}
\item $x$ est fermé vers le bas
\item $x$ est sans conflit
\end{itemize}
\pause

Pour la S.E suivante

\tikzfig{ex1}

Voici quelques configurations :
\begin{itemize}
\item $\{1,3\}$
\pause
\item $\{1,2,4,5\}$
\pause
\item $\{1,2,3,6\}$
\end{itemize}

\end{frame}

\begin{frame}
\frametitle{Jeu :}

$A$ une S.E munie d'une fonction

$\text{pol}_A : |A| \rightarrow \{-, +\}$
\pause

Par exemple sur l'exemple précédent

\tikzfig{ex1jeu}

est un jeu qu'on appelera $J_0$
\end{frame}

\begin{frame}
\frametitle{Partie :}

Une partie sur $A$ est un mot $s = s_1 \dots s_n \in |A|^*$ qui est :
\begin{itemize}
\item valide (à chaque étape on a une configuration)

\item non répétitive (on joue pas deux fois un événement)

\item alternante

\item négative (opposant commence)
\end{itemize}
\pause

\tikzfig{ex1jeu}

voici quelques parties
\begin{itemize}
\item $1236$
\pause
\item $1245$
\end{itemize}
\end{frame}

\begin{frame}
\frametitle{Construction sur les jeux}

Voici un autre jeu :
\tikzfig{ex2}

On l'appelera $(\mathbb{B} \rightarrow \mathbb{B})$


\end{frame}

\begin{frame}
\frametitle{Construction sur les jeux : jeu tenseur}

Voici le jeu $J_0 \otimes (\mathbb{B} \rightarrow \mathbb{B})$

\tikzfig{ex1jeu} \tikzfig{ex2}
\pause

Une partie sur ce jeu pourrait être $1 2 q_1 q_2 3 6$

On peut la restreindre à un des jeux

$1 2 q_1 q_2 3 6 \upharpoonright J_0 = 1 2 3 6$
\end{frame}


\begin{frame}
\frametitle{Construction sur les jeux : jeu dual}

Voici le jeu $J_0^\bot$

\tikzfig{ex1dual}

\end{frame}

\begin{frame}
\frametitle{Construction sur les jeux : jeu thèse}

On définit $A\vdash B := A^\bot \otimes B$

Voici le jeu $J_0\vdash (\mathbb{B} \rightarrow \mathbb{B})$

\tikzfig{ex1dual} \tikzfig{ex2}
\end{frame}

\begin{frame}
\frametitle{Exemples de stratégies sur $J_0\vdash  (\mathbb{B} \rightarrow \mathbb{B})$}

\tikzfig{ex1dual} \tikzfig{ex2}

Voici quelques stratégies :
\begin{itemize}
\item $\{\epsilon, q_1 1, q_1 1 2 3, q_1 1 2 3 6 q_2, q_1 4 5 tt_1\}$
\pause
\item $\{\epsilon,q_1 1, q_1 1 2 4\}$
\end{itemize}

\end{frame}

\begin{frame}
\frametitle{Stratégies}


Une stratégie sur le jeu $A$ est $\sigma \subset \mathcal{P}(A)$ qui est
\begin{itemize}
\item non vide
\item clos par préfixe pair
\item déterministe :
\begin{equation*} sa_1^+, sa_2^+ \in \sigma \implies a_1 = a_2 \end{equation*}
Le joueur a au plus un choix de coup pour rester dans la stratégie

\end{itemize}

\end{frame}

\begin{frame}
\frametitle{Catégories :}


Une catégorie $\mathcal{C}$ est la donnée de
\begin{itemize}
\item Une classe d'objets $|\mathcal{C}|$
\pause
\item Pour chaque paire d'objets $A, B \in |\mathcal{C}|$, une classe de
morphismes $\mathcal{C}(A,B)$
\pause

\item Pour chaque $A \in |\mathcal{C}|$, un morphisme
$\text{id}_A \in \mathcal{C}(A,A)$
\pause

\item Une loi $\circ$ sur les morphismes\\
$f \in \mathcal{C}(A,B)$, $g \in \mathcal{C}(B,C)$, alors
$g \circ f \in \mathcal{C}(A,C)$\\
telle que
\pause
\begin{itemize}
\item $\circ$ est associative:
\pause
\item Pour tout A, $\text{id}_A$ est neutre pour $\circ$
\end{itemize}
\end{itemize}
\end{frame}

\begin{frame}
\frametitle{Exemple de catégorie des jeux et stratégies}

Les objets sont les jeux

Un morphisme de $A$ vers $B$ est une stratégie sur le jeu $A\vdash B$
\pause
\tikzfig{excat}

\pause
Nous verrons ce qu'est la composition de stratégies et la stratégie $cc_A$ pour un jeu $A \vdash A$.
\end{frame}

\begin{frame}
\frametitle{Composition de stratégies :}


$\sigma$ stratégie sur $A \vdash B$, $\tau$ stratégie sur $B \vdash C$

$\sigma \circ \tau$ stratégie sur $A \vdash C$

\tikzfig{excompose}

\end{frame}

\begin{frame}
\frametitle{Composition de stratégies : Les interactions}

Une interaction sur $A, B, C$ est donnée par une séquence
$u \in (A \otimes B \otimes C)^*$ telle que
\begin{itemize}
\pause
\item $u \upharpoonright (A, B)$ est une partie sur $(A \vdash B)$
\item $u \upharpoonright (B, C)$ est une partie sur $(B \vdash C)$
\item $u \upharpoonright (A, C)$ est une partie sur $(A \vdash C)$
\end{itemize}

$I(A,B,C)$ est l'ensemble des interactions sur A,B,C.
\pause

Sur l'exemple précédent, on a l'interaction $123\dots 789$
et on a bien
\begin{itemize}
\item $12 \dots 89$ est une partie sur $A \vdash B$
\item $23 \dots 78$ est une partie sur $B \vdash C$
\item $13 \dots 79$ est une partie sur $A \vdash C$
\end{itemize}

\end{frame}

\begin{frame}
\frametitle{Composition de stratégies}


$\sigma \mid\mid \tau$ : interactions qui respectent $\sigma$ et $\tau$\\
($u \upharpoonright (A,B) \in \sigma$ et $u \upharpoonright (B,C) \in \tau$)
\pause

Alors la stratégie $\sigma \circ \tau$ est définie comme
\begin{equation*}\{u \upharpoonright (A, C) \mid u \in \sigma \mid\mid \tau \}\end{equation*}
\pause

Dans l'exemple précédent, $123\dots 789 \in \sigma \mid\mid \tau$, donc $13\dots79 \in \sigma \circ \tau$

\end{frame}


\begin{frame}
\frametitle{Stratégie copycat}

$cc_A$ stratégie sur $A\vdash A$ neutre pour la composition

\tikzfig{ex1dual} \tikzfig{ex1jeu}

\end{frame}

\begin{frame}
\frametitle{Définition inductive de parties : Jeu résiduel}


Le jeu $J_0$ :

\tikzfig{ex1jeu}

Le jeu $J_0\backslash 4$ :

\tikzfig{ex1resid}

On enlève 4 et les événements en conflit avec 4.
\end{frame}

\begin{frame}
\frametitle{Définition inductive de parties}

Parties commençant par un coup négatif : OPlay\\
Parties commençant par un coup positif : PPlay.

$\frac{}{\epsilon : \text{OPlay}_G}\epsilon_0$

$\frac{\text{min}(a^-,G) \qquad s : \text{PPlay}_{G\backslash a}}
{a^-s : \text{OPlay}_G}\text{OPlay}_i$

$\frac{\text{min}(a^+,G) \qquad s : \text{OPlay}_{G\backslash a}}
{a^+s : \text{PPlay}_G}\text{PPlay}_i$

\end{frame}

\begin{frame}
\frametitle{Exemple de partie}


Pour construire $1236 : \text{OPlay}_{J_0}$

\only<1->{
\begin{prooftree}
\hypo{}
\infer1{
    \epsilon : \text{OPlay}_{J_0\backslash 1,2,3,6}
}
\only<2->{
\hypo{\text{min}(6,J_0 \backslash 1,2,3)}
\infer2{
    6 : \text{PPlay}_{J_0 \backslash 1,2,3}
}
\only<3->{
\hypo{\text{min}(3,J_0 \backslash 1,2)}
\infer2{
    36 : \text{OPlay}_{J_0 \backslash 1,2}
}
}
}
\end{prooftree}
}
\only<3->{
\begin{prooftree}
\hypo{36 : \text{OPlay}_{J_0 \backslash 1,2}}
\only<4->{
\hypo{\text{min}(2,J_0 \backslash 1)}
\infer2{
    236:\text{PPlay}_{J_0 \backslash 1}
}
\only<5->{
\hypo{\qquad \text{min}(1,J_0)}
\infer2{
    1236 : \text{OPlay}_{J_0}
}
}
}
\end{prooftree}
}
\tikzfig{ex1proof}

\end{frame}

\begin{frame}
\frametitle{Cohérence entre parties (déterminisme de stratégie)}

$\frac{s:\text{Oplay}}{\text{Coh}^-(\epsilon,s)} \text{Coh}^-_{\epsilon,l} \qquad
 \frac{s:\text{OPlay}}{\text{Coh}^-(s,\epsilon)} \text{Coh}^-_{\epsilon,r}$

$\frac{a^- \neq a'^- \qquad s:\text{PPlay} \qquad s':\text{PPlay}}
{\text{Coh}^-(a^-s, a'^- s')} \text{Coh}^-_{\neq i}$

$\frac{s : \text{OPlay} \qquad s':\text{OPlay} \qquad \text{Coh}^-(s,s')}
{\text{Coh}^+(a^+s, a^+s')} \text{Coh}^+_i$

$\frac{s : \text{PPlay} \qquad s':\text{PPlay} \qquad \text{Coh}^+(s,s')}
{\text{Coh}^-(a^-s, a^-s')} \text{Coh}^-_i$
\pause

On prouve que ces deux propriétés sont équivalentes :
\begin{itemize}
\item Clos par préfixe et $\forall s, s' \in \sigma, \text{Coh}^-(s, s')$
\item Clos par préfixe et déterministe
\end{itemize}

\end{frame}

\begin{frame}
\frametitle{Définition inductive d'interactions}

Une interaction, un état (OOO,POP ou OPP), $u\upharpoonright (A,B), u \upharpoonright (B,C), u \upharpoonright (A,C)$
\resizebox{\textwidth}{!}{
\tikzfig{ex2} \tikzfig{ex1jeu} \tikzfig{ex1jeu}
}

\resizebox{0.5 \textheight}{!}{\tikzfig{diagpolbis}}
Pour $4_d 1_m q_1 tt_1 2_m 5_d$

\only<2>{$\epsilon : OOO$}
\only<3>{$5_d : OPP$}
\only<4>{$2_m 5_d : POP$}
\only<5>{$tt_1 2_m 5_d : OOO$}
\only<6>{$q_1 tt_1 2_m 5_d : POP$}
\only<7>{$1_m q_1 tt_1 2_m 5_d : OPP$}
\only<8>{$4_d 1_m q_1 tt_1 2_m 5_d : OOO$}
\end{frame}

\begin{frame}
\frametitle{Définition inductive d'interactions}



$\frac{}{\epsilon : OOOInt} OOO_{\epsilon}$

$\frac{s : POPInt \qquad a^+ : A}{a^+ s : OOOInt} aPOP \qquad
 \frac{s : OPPInt \qquad c^- : C}{c^- s : OOOInt} cOPP$

$\frac{s : OPPInt \qquad b^+ : B}{b^+ s : POPInt} bOPP \qquad
 \frac{s : OOOInt \qquad a^- : A}{a^- s : POPInt} aOOO$

$\frac{s : OOOInt \qquad c^+ : C}{c^+ s : OPPInt} cOOO \qquad
 \frac{s : POPInt \qquad b^- : B}{b^- s : OPPInt} bPOP$

\end{frame}

\begin{frame}
\frametitle{Stratégies résiduelles}

Pour la stratégie $\sigma = \{\epsilon, q_1 1, q_1 1 2 6, q_1 1 3 q_2\}$ sur $J_0 \vdash (\mathbb{B} \rightarrow \mathbb{B})$

\resizebox{0.5\textheight}{!}{
\tikzfig{ex1dual} \tikzfig{ex2}
}

On a $\sigma \backslash q_1 = \{1, 1 2 6, 1 3 q_2\}$ sur $J_0 \vdash (\mathbb{B} \rightarrow \mathbb{B})\backslash q_1$

\resizebox{0.5\textheight}{!}{
\tikzfig{ex1dual} \tikzfig{ex2resq1}
}

$s \in \sigma \backslash a \iff as \in \sigma$
\pause
On construit des OStrategy et PStrategy
\end{frame}

\begin{frame}
\frametitle{Conclusion : Ce qui a été fait}

Construire la stratégie composée
\pause

Preuve que c'est bien une stratégie.
\pause

Pour les preuves, principe d'induction sur les structures inductives mutuelles.
\end{frame}

\begin{frame}
\frametitle{Conclusion : Ce qui reste à faire}

\begin{itemize}
\item Sanity Check
\pause
\item Associativité de la composition
\pause
\item Construction de Copycat
\pause
\item Preuve que Copycat est neutre
\end{itemize}
\pause
Définitions inductives efficaces pour les preuves, plein d'espoir.
\end{frame}

\begin{frame}
\frametitle{FIN}

Merci d'avoir écouté
\end{frame}

\end{document}
